\section{Метод решения}
Робот --- скрипт \texttt{crawler.py} на Python~3 (стандартная библиотека, PyYAML и драйвер MongoDB \texttt{pymongo}), запуск: \texttt{python3 crawler.py config.yaml}. База данных --- распределённая: MongoDB, набор реплик из трёх узлов (раздел \enquote{База данных}). Каждый источник обходится в ширину отдельным потоком со своей задержкой между запросами, так что источники качаются параллельно, а к одному сайту одновременно идёт не больше одного запроса.

\subsection{Конфиг}
\begin{itemize}
    \item \texttt{db.uri} --- адрес набора реплик MongoDB (все три узла и имя набора \texttt{rs0}); \texttt{db.name} --- имя базы.
    \item \texttt{logic}: \texttt{delay} --- пауза между запросами к источнику (0,5~с); \texttt{recrawl\_after} --- через сколько секунд после последней проверки документ проверяется снова (7 дней); \texttt{timeout}, \texttt{user\_agent}.
    \item \texttt{sources} --- список источников: \texttt{name}; \texttt{seeds} --- начальные URL; \texttt{follow} --- регулярное выражение для страниц, с которых берутся ссылки; \texttt{save} --- для страниц, сохраняемых как документы; \texttt{links\_in} --- \texttt{id} элемента, внутри которого искать ссылки; \texttt{max\_depth} --- максимальная глубина обхода; \texttt{limit} --- сколько документов собрать.
\end{itemize}

Источники пробного прогона (тематика --- туризм, как в ЛР1):
\begin{enumerate}
    \item Wikivoyage: начало --- \texttt{/wiki/Destinations}, обходятся и сохраняются все статьи \texttt{https://en.wikivoyage.org/wiki/<название>} без двоеточия в названии (служебные пространства имён \texttt{Special:}, \texttt{File:}, \texttt{Wikivoyage:} и т.\,п. отсекаются).
    \item Wikipedia: начало --- \texttt{/wiki/Category:Tourism}; ссылки берутся только со страниц категорий (подкатегории и статьи), сохраняются статьи. Глубина до 4: категория, три уровня подкатегорий и статьи из них.
\end{enumerate}
Ссылки берутся только внутри \texttt{<div id="mw-content-text">}, то есть из текста статьи или списка категории, без меню сайта.

\subsection{База данных}
Распределённая база --- MongoDB~7.0, набор реплик \texttt{rs0} из трёх узлов: по контейнеру Docker на узел, порты 27017--27019 (\texttt{docker-compose.yml}, инициализация набора --- \texttt{scripts/mongo\_init.py}). Один узел основной (PRIMARY) и принимает запись, два вторичных (SECONDARY) держат полные копии данных и догоняют основной по журналу операций. Если основной узел пропадает, оставшиеся выбирают новый (около 10~с), а драйвер сам переключается на него и повторяет прерванные операции. Робот пишет с подтверждением большинства (\texttt{w=majority}): запись считается выполненной, когда она есть хотя бы на двух узлах из трёх, поэтому отказ одного узла её не теряет.

Коллекции (ключ \texttt{\_id} --- нормализованный URL):
\begin{itemize}
    \item \texttt{documents}: \texttt{url}, \texttt{source}, \texttt{html} (сырой HTML страницы), \texttt{fetched\_at} (дата обкачки, Unix time), \texttt{checked\_at} (дата последней проверки на изменение), \texttt{last\_modified} и \texttt{etag} (заголовки ответа для условного запроса), \texttt{hash} (SHA-1 HTML). Индекс \texttt{(source, checked\_at)} --- для поиска документов, которые пора перепроверить.
    \item \texttt{frontier} --- очередь обхода: \texttt{source}, \texttt{depth}, \texttt{seq} --- номер постановки в очередь (растёт, берётся из времени в наносекундах), \texttt{status}: \texttt{new} --- ждёт обкачки, \texttt{done} --- обработан, \texttt{dup} --- дубликат уже сохранённого документа, \texttt{error} --- ошибка HTTP. Порядок обхода --- по \texttt{(depth, seq)}, то есть в ширину; индекс \texttt{(source, status, depth, seq)}.
\end{itemize}
Новая ссылка добавляется операцией upsert с \texttt{\$setOnInsert}: если URL уже в очереди, запись не меняется.

\subsection{Шаг обхода}
\begin{enumerate}
    \item Взять следующий URL: сначала документ, который пора перепроверить, затем первый URL очереди со статусом \texttt{new}.
    \item Скачать страницу (\texttt{GET} с \texttt{User-Agent} и \texttt{Accept-Encoding: gzip}). При сетевой ошибке, ответе 429 или 5xx --- до 3 попыток с паузой. Страница не \texttt{text/html} или другой код ответа --- статус \texttt{error}.
    \item Определить кодировку: \texttt{charset} из \texttt{Content-Type}, затем \texttt{<meta charset>}, иначе UTF-8.
    \item Разобрать HTML (\texttt{html.parser}): ссылки \texttt{<a href>} и \texttt{<link rel="canonical">}. Перенаправления MediaWiki отдают ту же статью под старым URL с ответом 200, поэтому документ сохраняется под каноническим URL. Если такой уже есть в базе, страница помечается \texttt{dup}.
    \item Ссылки нормализуются и попадают в очередь, если подходят под \texttt{follow} или \texttt{save}, разрешены \texttt{robots.txt} и не превышают \texttt{max\_depth}. Дубликаты отсекает уникальный ключ (\texttt{INSERT OR IGNORE}).
    \item Документ, новые ссылки и статус URL записываются в одной многодокументной транзакции MongoDB (она доступна именно в наборе реплик); при конфликте записи транзакция повторяется. Затем пауза \texttt{delay}.
\end{enumerate}

\textbf{Нормализация URL:} относительные ссылки разрешаются от URL страницы (\texttt{urljoin}, в т.\,ч. \texttt{./}, \texttt{../}, \texttt{//host}); схема и хост приводятся к нижнему регистру; удаляются порт по умолчанию и якорь \texttt{\#\ldots}; путь перекодируется в единый вид (\texttt{S\%c3\%a3o}, \texttt{São} $\to$ \texttt{S\%C3\%A3o}, \texttt{\%27} $\to$ \texttt{'}); параметры запроса сортируются; пустой путь заменяется на \texttt{/}; ссылки не \texttt{http(s)} (\texttt{mailto:}, \texttt{javascript:}) отбрасываются.

\textbf{Вежливость:} \texttt{robots.txt} каждого хоста скачивается один раз за запуск и проверяется \texttt{urllib.robotparser}. Если в нём задан \texttt{Crawl-delay} больше \texttt{delay}, берётся он.

\subsection{Остановка и продолжение}
Очередь и статусы хранятся в базе, страница обрабатывается одной транзакцией. По Ctrl+C или SIGTERM робот дожидается конца текущего запроса и завершается. При аварийной остановке (\texttt{kill -9}) незавершённая транзакция откатывается, URL остаётся в статусе \texttt{new}. Повторный запуск продолжает обход со следующего необработанного URL, а уже сохранённые страницы не качает.

\subsection{Переобкачка}
Документ считается устаревшим, если с последней проверки прошло больше \texttt{recrawl\_after} секунд. Такие документы проверяются раньше новых URL условным запросом \texttt{GET} с \texttt{If-Modified-Since} (значение \texttt{Last-Modified} прошлого ответа) и \texttt{If-None-Match} (если был \texttt{ETag}):
\begin{itemize}
    \item 304 Not Modified --- документ не изменился, обновляется только \texttt{checked\_at};
    \item 200 --- сравнивается SHA-1 нового HTML с сохранённым; если отличается, обновляются \texttt{html}, \texttt{fetched\_at} и заголовки, иначе только \texttt{checked\_at}.
\end{itemize}
За один запуск документ проверяется не больше одного раза, поэтому запуск всегда завершается. Периодичность обеспечивается повторными запусками (например, из \texttt{cron}): каждый запуск перепроверяет документы старше \texttt{recrawl\_after} и докачивает очередь.

\pagebreak
